\par\Needspace{6\baselineskip}\medskip
\begin{center}\hypertarget{native-196-article146}{}\pdfbookmark[2]{146.}{article-146}146.\end{center}
Calculus elementorum, hinc ex $r',r'',2f$ atque temporis intervallo correcto
inter observationem secundam et tertiam, cuius productum in quantitatem $k$
(art. 1) per $\theta$ denotamus, illinc ex $r,r',2f''$ atque temporis intervallo inter observationem
primam et secundam, cuius productum per $k$ esto $=\theta''$, secundum
methodum in artt.$^{\prime}$ 88—105 expositam tantummodo usque ad quantitatem illic
per $y$ denotatam producendus est, cuius valorem in combinatione priori per $\eta$,
in posteriori per $\eta''$ denotabimus. Fiat deinde
\[
\frac{\theta''\eta}{\theta\eta''}=P',\qquad
\frac{r'r'\theta\theta''}{rr''\eta\eta''\cos f\cos f'\cos f''}=Q',
\]
\SourceLeaf{197}patetque, si valores quantitatum $P,Q$, quibus totus hucusque calculus superstructus
erat, ipsi veri fuerint, evadere debere $P'=P,Q'=Q$. Vice versa
facile perspicitur, si prodeat $P'=P,Q'=Q$, duplicem elementorum calculum,
si utrimque ad finem perducatur, numeros prorsus aequales suppeditaturum
esse, per quos itaque omnes tres observationes exacte repraesentabuntur, adeoque
problemati ex asse satisfiet. Quoties autem non fit $P'=P,Q'=Q$, accipientur
$P'-P,Q'-Q$ pro $X$ et $Y$, siquidem $P$ et $Q$ pro $x$ et $y$ acceptae fuerint:
adhuc magis commodum erit statuere $\log P=x$, $\log Q=y$, $\log P'-\log P=X$, $\log Q'-\log Q=Y$.
Dein calculus cum aliis valoribus ipsarum $x,y$ repetendus erit.
